\documentclass[11pt, a4paper]{article}

\usepackage{fontspec}
\setmainfont{Latin Modern Roman}
\setmonofont{DejaVu Sans Mono}[Scale=0.85]
\usepackage{polyglossia}
\setdefaultlanguage{french}
\setotherlanguage{english}
\usepackage[margin=2.5cm]{geometry}
\usepackage{amsmath, amssymb, amsthm}
\usepackage{xurl}
\usepackage{hyperref}
\usepackage[dvipsnames]{xcolor}
\usepackage{enumitem}
\usepackage{fancyvrb}
\usepackage{caption}
\usepackage{draftwatermark}
\SetWatermarkText{PRÉPUBLICATION}
\SetWatermarkScale{1.0}
\SetWatermarkLightness{0.92}
\captionsetup[figure]{name=Listing}
\emergencystretch=5em
\newcommand{\brk}{\discretionary{}{}{}}
\newcommand{\lean}[1]{\ifmmode\text{\ttfamily\detokenize{#1}}\else{\ttfamily\detokenize{#1}}\fi}

\definecolor{indigo}{rgb}{0.29, 0.0, 0.51}
\definecolor{teal}{rgb}{0.0, 0.5, 0.5}
\hypersetup{
    colorlinks=true,
    linkcolor=teal,
    citecolor=indigo,
    urlcolor=blue,
    pdftitle={Quotients eta Fricke-auto-duaux (prepublication)},
    pdfauthor={Xavier Callens}
}

\newtheorem{theorem}{Théorème}[section]
\newtheorem{lemma}[theorem]{Lemme}
\newtheorem{definition}[theorem]{Définition}
\newtheorem{corollary}[theorem]{Corollaire}
\theoremstyle{remark}
\newtheorem{remark}[theorem]{Remarque}
\renewcommand{\proofname}{Démonstration}

\title{\textbf{Les quotients êta Fricke-auto-duaux sont des formes propres}\\[0.3em]
  \large\textbf{Un corollaire en Lean 4, et une note sur sa motivation \\ dans la dualité S de
  Persson--Volpato pour les modèles CHL}\\[0.8em]
  \large\textcolor{BrickRed}{PRÉPUBLICATION --- non relue par les pairs}}
\author{
  Xavier Callens\thanks{Auteur correspondant : \texttt{callensxavier@gmail.com}} \\
  \small SocrateAI Research Initiative
}
\date{\today\\[0.6em]
  \small Archivé sur Zenodo : \texttt{ZENODO\_DOI\_PLACEHOLDER} (dépôt en cours)}

\begin{document}
\maketitle

\begin{center}
\fbox{\begin{minipage}{0.92\textwidth}
\small \textbf{Statut.} Ceci est une \emph{prépublication}. Elle n'a pas fait l'objet d'une
\textbf{relecture par les pairs}. Le développement Lean qu'elle décrit compile et son empreinte
axiomatique est auditée --- ces affirmations sont vérifiées par machine et reproductibles de façon
indépendante --- mais l'\emph{exposition} n'a reçu aucune relecture externe. Le contenu
mathématique est un unique corollaire clos (Théorème~\ref{thm:main}) d'une loi de transformation
déjà démontrée ; la motivation physique citée au \S\ref{sec:motivation} est rapportée honnêtement
comme une \emph{motivation}, non comme quelque chose que cet article formalise ou démontre. Lire
le périmètre avant de citer le résultat.
\end{minipage}}
\end{center}

\begin{abstract}
Nous formalisons, en Lean~4 au-dessus de Mathlib, un corollaire d'une loi de transformation de
Fricke déjà démontrée pour les quotients êta : si le vecteur d'exposants $r$ du quotient êta
$f(z) = \prod_{\delta\mid N}\eta(\delta z)^{r_\delta}$ vérifie $r_\delta = r_{N/\delta}$ pour tout
$\delta\mid N$ (\emph{Fricke-auto-dualité}), alors $f$ est une véritable forme propre de la
transformation de Fricke $z\mapsto -1/(Nz)$, avec une valeur propre en forme close
$\lambda = i^{-k}N^{k/2}$ ne dépendant que du niveau $N$ et du poids $k$, non de $r$. Nous
démontrons également précisément quand $\lambda=\pm1$. Le seul contenu mathématique nouveau est
une unique identité de produit, $\prod_{\delta\mid N}\delta^{r_\delta}=N^k$ sur les $r$
auto-duaux ; toute autre déclaration est une substitution directe dans la loi de transformation
préexistante. Nous rapportons, comme motivation et au niveau littérature de l'échelle
épistémique de ce programme --- non comme une affirmation formalisée --- que cette condition
d'auto-dualité est exactement le « Frame shape équilibré » de Persson et Volpato, qu'ils
identifient comme caractérisant les compactifications de cordes de type CHL auto-duales sous une
S-dualité de type Fricke sur le module axio-dilaton hétérotique. Nous énonçons précisément où
leur normalisation et la nôtre divergent, afin que les deux ne soient pas confondues.
\end{abstract}

\noindent\textbf{MSC 2020 :} 11F11, 11F20, 68V20.\quad
\textbf{Mots-clés :} quotients êta, involution de Fricke, formes propres, Lean 4, Mathlib,
vérification formelle, modèles CHL.

\section{Introduction}\label{sec:introduction}

Un développement compagnon de ce programme, \emph{Une formalisation idiomatique en Mathlib de
l'involution de Fricke}~\cite{Fricke2026}, formalise l'involution de Fricke $W_N$ au niveau
matriciel sur $\Gamma_0(N)$, et un second, \emph{Les quotients êta en Lean 4}~\cite{EtaQuotients2026},
formalise la loi de transformation des quotients êta sous $W_N$ : pour
$f(z)=\prod_{\delta\mid N}\eta(\delta z)^{r_\delta}$ de poids $k$ (c.-à-d.
$\sum_{\delta\mid N} r_\delta = 2k$),
\begin{equation}\label{eq:etaQuotient_fricke}
  f\!\left(\frac{-1}{Nz}\right) \;=\; i^{-k}N^k\Bigl(\sqrt{\textstyle\prod_{\delta\mid N}\delta^{r_\delta}}\Bigr)^{-1} z^k \, f^\ast(z),
\end{equation}
où $f^\ast$ est le quotient \emph{dual}, de vecteur d'exposants $\delta\mapsto r_{N/\delta}$.
C'est \lean{etaQuotient_fricke} (étiquette \texttt{F3.2-A7}), déjà démontré, sans \lean{sorry},
dans \texttt{EtaQuotientModularity.lean}. En général $f^\ast\neq f$ : \lean{fricke_dual_pin} dans
ce module exhibe un $r$ dont le dual est une fonction authentiquement différente.

Cette note pose la question naturelle suivante : que se passe-t-il quand $f^\ast=f$ ?
L'équation~\eqref{eq:etaQuotient_fricke} s'effondre alors en une véritable relation de forme
propre, avec une valeur propre ne dépendant que de $N$ et $k$. Cet effondrement est court --- une
substitution et une simplification, non une dérivation nouvelle --- et son seul ingrédient non
trivial est une unique identité de produit (Lemme~\ref{lem:prod-collapse}). Nous le formalisons,
avec une caractérisation close de quand la valeur propre vaut exactement $\pm1$, et nous
consignons, honnêtement et à distance, pourquoi la question valait la peine d'être posée : la
condition d'auto-dualité sur $r$ est exactement la forme combinatoire d'une condition de la
littérature physique que cet article ne cherche pas à dériver.

\section{La condition d'auto-dualité}

\begin{definition}[Vecteur d'exposants Fricke-auto-dual]\label{def:selfdual}
\lean{IsFrickeSelfDual}. Pour $N>0$ et un vecteur d'exposants $r:\mathbb{N}\to\mathbb{Z}$,
\[
  \mathrm{IsFrickeSelfDual}(N,r) \;:\Longleftrightarrow\; \forall\, \delta\mid N,\ r_\delta = r_{N/\delta}.
\]
\end{definition}

Nous appelons cette condition \emph{auto-duale} (de façon équivalente \emph{Fricke-symétrique}),
jamais « équilibrée ». \texttt{EtaQuotientModularity.lean} utilise déjà « équilibrée »
(\emph{balanced}) pour un concept sans rapport --- \lean{exists_balanced_add}, étiquette
\texttt{F3.2-B1}, l'étape de division à reste équilibré (\lean{Int.bmod}) utilisée dans la
descente des générateurs de $\Gamma_0(N)$ --- et la littérature physique discutée au
\S\ref{sec:motivation} utilise « équilibré » (\emph{balanced}) dans un troisième sens, apparenté
mais distinct. Réutiliser le mot ici entrerait en collision avec les deux. \lean{IsFrickeSelfDual}
est \lean{Decidable} (\lean{decidableIsFrickeSelfDual}), si bien que l'appartenance peut être
vérifiée par \lean{decide} sur des instances concrètes, ce que fait chaque épingle
(\emph{pin}) du développement ci-dessous.

\begin{lemma}[Le produit s'effondre en une puissance du niveau]\label{lem:prod-collapse}
\lean{prod_zpow_selfDual}. Pour $N>0$, $r$ auto-dual sur $N$, et $\sum_{\delta\mid N} r_\delta=2k$,
\[
  \prod_{\delta\mid N} \delta^{r_\delta} \;=\; N^k.
\]
\end{lemma}

C'est le seul endroit où l'hypothèse d'auto-dualité effectue un véritable travail au-delà d'une
substitution : elle force le produit des exposants pondéré par les diviseurs à s'effondrer en une
unique puissance du niveau, ce qui est ce qui retire $r$ de la valeur propre entièrement. Ceci a
été vérifié, indépendamment de la preuve Lean, en arithmétique rationnelle exacte sur quatre
plages (jusqu'à $1\le N\le 60$, exposants dans $[-4,4]$, $671504$ instances auto-duales), zéro
désaccord, et contre $9597$ vecteurs de poids pair non auto-duaux dont $98{,}8\%$ violent
l'identité --- confirmant que l'auto-dualité porte réellement l'effondrement, et n'est pas un
artefact de la seule condition de poids.

\section{La relation de forme propre}

\begin{definition}[Valeur propre de Fricke]\label{def:eigenvalue}
\lean{frickeEigenvalue}. Pour $N>0$ et $k\in\mathbb{Z}$,
\[
  \lambda(N,k) \;:=\; i^{-k}\sqrt{N^{k}}.
\]
\end{definition}

\begin{theorem}[Les quotients êta auto-duaux sont des formes propres de Fricke]\label{thm:main}
\lean{etaQuotient_fricke_selfDual}. Pour $N>0$, $r$ auto-dual sur $N$, $\sum_{\delta\mid N}
r_\delta=2k$, et $z$ dans le demi-plan de Poincaré supérieur ouvert,
\[
  f\!\left(\frac{-1}{Nz}\right) \;=\; \lambda(N,k)\, z^k\, f(z), \qquad f(z)=\prod_{\delta\mid N}\eta(\delta z)^{r_\delta}.
\]
\end{theorem}

La démonstration tient en deux réécritures : substituer, dans l'équation~\eqref{eq:etaQuotient_fricke},
le quotient dual $f^\ast$ par $f$ lui-même (la Définition~\ref{def:selfdual} fait de ceci une
congruence sur le produit indexé par les diviseurs, \lean{etaQuotient_congr_divisors}), puis
effondrer la constante à l'aide du Lemme~\ref{lem:prod-collapse} (\lean{fricke_const_selfDual}).
Rien d'autre n'intervient. Une forme intermédiaire non normalisée, conservant
$\sqrt{\prod\delta^{r_\delta}}$ plutôt que $\lambda$, est \lean{etaQuotient_fricke_selfDual_raw} ;
une forme normalisée isolant la racine de l'unité $i^{-k}$ d'un côté est
\lean{etaQuotient_fricke_selfDual_normalized}.

\begin{remark}[Ce que ceci n'est pas]\label{rem:not-modularity}
Le Théorème~\ref{thm:main} ne fait aucune référence à $\Gamma_0(N)$, aux conditions aux pointes,
ni au type empaqueté $M_k(\Gamma_0(N))$ : c'est une affirmation portant sur la valeur d'une
fonction en un point, $-1/(Nz)$, transportée depuis la loi de transformation générale déjà
démontrée. Appeler $\lambda(N,k)$ « la valeur propre de l'opérateur de Fricke $W_N$ sur
$M_k(\Gamma_0(N))$ » exagérerait ce qui est démontré ici ; la lecture correcte est que $f$ est un
vecteur propre de la seule application linéaire $g\mapsto g(-1/(Nz))$ restreinte aux fonctions de
cette forme.
\end{remark}

\subsection{Quand la valeur propre vaut-elle $\pm1$ ?}

La dépendance de la valeur propre en $N$ et $k$ seuls rend accessible une caractérisation exacte
de $\lambda=\pm1$ sans hypothèse supplémentaire sur $r$.

\begin{theorem}[Module, et les critères de signe]\label{thm:pm1}
Pour $N>0$ et $k\in\mathbb{Z}$ :
\begin{align}
\|\lambda(N,k)\| &= \sqrt{N^{k}} \tag{\lean{frickeEigenvalue_norm}}\\
\lambda(N,k)=1 &\iff 4\mid k \ \wedge\ (N=1 \vee k=0) \tag{\lean{frickeEigenvalue_eq_one_iff}}\\
\lambda(N,k)=-1 &\iff k\bmod 4 = 2 \ \wedge\ N=1 \tag{\lean{frickeEigenvalue_eq_neg_one_iff}}
\end{align}
\end{theorem}

\begin{corollary}[$\pm1$ en fonction de la somme des poids]\label{cor:pm1-weight}
\lean{frickeEigenvalue_pm_one_iff}. Pour $N>0$, $r$ auto-dual, $\sum_{\delta\mid N} r_\delta=2k$,
\[
  \lambda(N,k)\in\{1,-1\} \iff 4 \Bigm| \textstyle\sum_{\delta\mid N} r_\delta \ \wedge\ \bigl(N=1 \vee \sum_{\delta\mid N} r_\delta = 0\bigr).
\]
\end{corollary}

Le Théorème~\ref{thm:pm1} repose sur deux faits arithmétiques indépendants, chacun démontré et
épinglé séparément avant d'être combiné : la périodicité de $i^{k}\bmod 4$
(\lean{I_zpow_emod}, authentiquement nouveau --- le propre \lean{Complex.I_pow_eq_pow_mod} de
Mathlib n'est énoncé que pour des exposants dans $\mathbb{N}$, non $\mathbb{Z}$), et
$(N:\mathbb{R})^k=1 \iff (N=1\vee k=0)$ pour $N>0$ (\lean{natCast_zpow_eq_one_iff}). Chaque
instance épinglée dans les deux théorèmes est vérifiée deux fois, par le lemme général et par un
calcul manuel indépendant ne partageant aucun lemme avec lui ; les instances informatives sont
celles où une supposition naïve serait fausse --- par exemple $(N,k)=(6,2)$, où $4\mid k$ échoue
mais $k\bmod 4=2$ tient et $\lambda\neq -1$ uniquement parce que $N\neq1$, le cas qui teste le
plus directement si la condition de niveau a été omise par accident.

\section{État de la vérification}\label{sec:verification}

Le développement compte 333 déclarations en 6437 lignes, dans un nouveau module,
\texttt{EtaQuotient\brk FrickeSelfDual.lean}, importé dans la cible de compilation par défaut de
la bibliothèque. Il est entièrement sans \lean{sorry} : 24 nœuds du graphe d'énoncés (deux
définitions, seize lemmes généraux et caractérisations, cinq épingles de valeurs propres
vérifiées croisées contre une évaluation en produit-$q$ de $\eta$, et une garde de forme
d'énoncé), tous \texttt{proved}, aucun bloqué, aucun ouvert. 299 des 333 déclarations portent une
garde \lean{\#guard\_msgs in \#print axioms} qui fait échouer la compilation, rapportant
exactement \lean{[propext, Classical.choice, Quot.sound]} ; le reste sont les définitions de
données pures (les vecteurs d'exposants épinglés) et des lemmes auxiliaires qu'une garde ne
renforcerait pas. \lean{lake build SocrateAI} --- la cible par défaut partagée de toute la
bibliothèque, portant désormais aussi les gardes de ce module --- s'achève en 3769 tâches sans
erreur.

Chaque théorème principal a été vérifié, indépendamment de sa preuve Lean, contre une évaluation
directe en produit-$q$ de la fonction êta de Dedekind (troncature à 600 termes, trois points par
instance, erreur relative $\le 1{,}8\times10^{-14}$), et l'instance $N=1$, $r\equiv24$, $k=12$ du
Théorème~\ref{thm:main} reproduit \lean{eta_S_via_fricke}, lui-même re-dérivé indépendamment du
propre \lean{discriminant_S_invariant} de Mathlib. Les termes de preuve du noyau ont été
inspectés directement (\lean{\#print}) pour confirmer que les théorèmes principaux sont
d'authentiques substitutions dans \lean{etaQuotient_fricke} plutôt que des reformulations qui
compilent simplement : chacun se réduit à au plus deux étapes \lean{rw} suivies de \lean{ring}.

Le graphe de dépendance des énoncés (\lean{dag/theorems.jsonl}) est valide : 104 nœuds, 95
démontrés, acyclique. Aucun axiome au-delà des trois axiomes standard de Lean et des six
justificatifs \lean{physics\_postulate\_$i$} préexistants et déclarés de la bibliothèque (inutilisés
par quoi que ce soit démontré ici) n'est déclaré nulle part dans la bibliothèque.

\section{Motivation : la S-dualité de Fricke de Persson--Volpato dans les modèles CHL}\label{sec:motivation}\label{sec:selfdual-pointer}

\textbf{Rien dans cette section n'est formalisé. Rien dans les
\S\S\ref{def:selfdual}--\ref{sec:verification} ne nomme ni ne dépend d'un objet introduit ici.}
Nous rapportons la motivation qui a conduit à poser la question, au niveau littérature (L) de
l'échelle épistémique de ce programme, car la norme de transparence que ce programme
suit~\cite{Tao2026AI,Leiden2026} exige d'attribuer les sources clairement plutôt que de laisser
un artefact formel suggérer une provenance qu'il n'a pas.

Persson et Volpato~\cite{PerssonVolpato2015} étudient les compactifications CHL --- des orbifolds
de la corde hétérotique (de façon équivalente, par dualité de cordes, du type~IIA sur
$K3\times T^2$) par une symétrie d'ordre $N$, $g$ --- et attachent à $g$ un « Frame shape
généralisé » $g\leftrightarrow\prod_{a\mid N}a^{m(a)}$, un vecteur d'exposants sur les diviseurs
de $N$, contraint par $\sum_{a\mid N}m(a)\cdot a=24$ (leur éq.~1.4, consignant une représentation
de dimension $24$). Ils appellent le Frame shape \emph{équilibré} (\emph{balanced}, leur éq.~1.6)
quand $m(N/a)=m(a)$ pour tout $a\mid N$ --- exactement la Définition~\ref{def:selfdual}, lue sur
leur vecteur d'exposants --- et leur résultat physique principal (leurs \S1.1, \S3) est qu'un
modèle CHL est auto-dual sous une S-dualité de Fricke $S\mapsto -1/(NS)$ sur le module
axio-dilaton hétérotique $S$ précisément quand son Frame shape est équilibré. Cette
identification est établie par des arguments physiques que cette bibliothèque Lean ne contient
pas et que cette note ne cherche pas à reproduire : $N$-modularité du réseau de charges, théorie
des représentations du groupe de Conway sur la représentation de dimension $24$, et comptage
d'états BPS / indice de Witten.

\begin{remark}[Deux normalisations à ne pas confondre]\label{rem:normalization}
La contrainte de poids de Persson--Volpato, $\sum_{a\mid N}m(a)\cdot a=24$, n'est \emph{pas} la
normalisation de poids des quotients êta, $\sum_{\delta\mid N}r_\delta=2k$, utilisée dans tout cet
article. Les deux ne coïncident qu'à $N=1$. Leur Frame shape équilibré canonique à $N=2$,
$1^8 2^8$, vérifie $\sum a\cdot m(a) = 8+16=24$ mais $\sum m(a)=16$, c.-à-d. un poids $8$, non
$12$ (vérifié cette session sur $1^k N^k$, $k=24/(1+N)$, pour $N\in\{1,2,3,5,7,11,23\}$). Un
lecteur traduisant entre les deux littératures doit reporter le poids réel du vecteur
d'exposants, non supposer un poids $12$ partout.
\end{remark}

Le module $S$ chez Persson--Volpato est l'axio-dilaton hétérotique, non le module de structure
complexe $\tau$ du $T^2$ qu'un résultat négatif compagnon de ce
programme~\cite{Fricke2026} a examiné et trouvé \emph{non} soutenu, au niveau $N>1$, par la
référence naturelle en dualité T (Giveon--Porrati--Rabinovici). Que la même matrice de Fricke
$W_N=\left(\begin{smallmatrix}0&-1\\N&0\end{smallmatrix}\right)$ revienne attachée à un module
\emph{différent} dans une référence correctement sourcée, tout en échouant à s'attacher au
module du $T^2$ dans une référence incorrectement sourcée, est exactement la distinction que ce
résultat négatif était conçu pour faire apparaître, et la discipline de citation de cette note en
est une continuation directe : citer la source qui soutient effectivement l'affirmation précise
motivée, non la source qui mentionne seulement une matrice d'apparence similaire.

\section{Travaux apparentés}\label{sec:related}

Martin~\cite{Martin1996} donne le compte rendu classique de la multiplicativité des valeurs
propres de Fricke et d'Atkin--Lehner pour les quotients êta ; le Théorème~\ref{thm:main} en est
le cas particulier directement accessible depuis la loi de transformation générale déjà
démontrée, énoncé et vérifié en Lean plutôt que re-dérivé. \lean{etaQuotient_fricke}
(\S\ref{sec:introduction})\cite{EtaQuotients2026} fournit la loi générale que cette note
spécialise ; \cite{Fricke2026} fournit l'involution de Fricke $W_N$ au niveau matriciel à laquelle
les lois de transformation des deux articles sur les quotients êta sont finalement rapportées, et
son propre résultat négatif sur la dualité T est le précédent direct de la discipline de citation
de cette note au \S\ref{sec:motivation}.

\section{Limites}\label{sec:limitations}

\begin{remark}[Ce qui n'est pas fait ici]
\leavevmode
\begin{enumerate}[label=(\alph*)]
  \item Aucun lien avec la modularité $\Gamma_0(N)$, le comportement aux pointes, ou le type
    empaqueté $M_k(\Gamma_0(N))$ n'est établi (Remarque~\ref{rem:not-modularity}).
  \item L'affirmation physique motivant cette note (\S\ref{sec:motivation}) n'est pas formalisée,
    n'est pas tentée ici, et ne doit pas être inférée du contenu Lean.
  \item $\lambda(0,k)=0$ est une valeur parasite (la définition est totale mais l'hypothèse $N>0$
    est requise par chaque théorème qui l'utilise) ; aucun théorème ici ne repose sur, ni
    n'affirme quoi que ce soit à propos de $N=0$.
  \item La caractérisation de $\lambda=\pm1$ (Théorème~\ref{thm:pm1}) est complète pour la valeur
    propre elle-même, mais le Corollaire~\ref{cor:pm1-weight} hérite de l'hypothèse
    d'auto-dualité sur $r$ --- il ne caractérise pas quels vecteurs d'exposants sont auto-duaux,
    une question combinatoire distincte que cette note n'aborde pas au-delà de la décidabilité de
    la Définition~\ref{def:selfdual}.
\end{enumerate}
\end{remark}

\section{Conclusion}

Nous avons formalisé, en Lean~4, qu'un quotient êta Fricke-auto-dual est une forme propre de la
transformation de Fricke, avec une valeur propre en forme close ne dépendant que du niveau et du
poids, et précisément quand cette valeur propre vaut $\pm1$. Les mathématiques sont un corollaire
court et honnête d'une loi de transformation déjà démontrée ; leur intérêt tient moins à la
difficulté de la preuve qu'au fait que la même condition combinatoire organise une dualité
physique dans la littérature des compactifications de cordes CHL, rapportée ici strictement comme
motivation, avec les normalisations divergentes des deux littératures énoncées explicitement afin
qu'elles ne puissent être confondues silencieusement.

\section*{Note sur la méthode : assistance de l'IA et responsabilité humaine}

Cette note suit la norme de transparence défendue par Tao~\cite{Tao2026AI}, qui reprend les
recommandations de la Déclaration de Leyde sur l'intelligence artificielle et les
mathématiques~\cite{Leiden2026} : divulguer l'usage d'outils, affirmer que le crédit et la
responsabilité reviennent à l'auteur humain, et attribuer les sources là où les outils d'IA ne
peuvent pas le faire de manière fiable eux-mêmes.

\textbf{Ce que l'IA a fait.} La formalisation en Lean~4, la recherche bibliographique et la
rédaction de ce texte ont été produites avec l'assistance de grands modèles de langage (la famille
Claude d'Anthropic), sous direction humaine sur plusieurs sessions. La méthodologie est
neuro-symbolique au sens précis que décrit Tao~\cite{Tao2026AI} : le raisonnement en langage
naturel propose des définitions, des stratégies de preuve et de la prose, et c'est le noyau Lean
--- non le modèle, ni la confiance de l'auteur dans le modèle --- qui admet ou rejette chaque
affirmation mathématique. Chaque théorème de cet article compile contre Mathlib sans \lean{sorry}
et avec une garde d'axiomes qui fait échouer la construction en cas de dérive
(\S\ref{sec:verification}) ; rien ici n'est affirmé vrai parce qu'un modèle l'a dit.

\textbf{Ce que l'auteur a fait.} Xavier Callens a fixé l'orientation de la recherche --- en
choisissant cette spécialisation précisément parce qu'elle était accessible depuis un outillage
déjà démontré plutôt que d'exiger une nouvelle tentative de pont physique, après qu'une exécution
compagne de ce programme eut trouvé une telle tentative correctement refusée~\cite{Fricke2026}
--- a jugé ce que signifie le résultat et s'il mérite d'être publié, et est responsable de chaque
affirmation de ce document, y compris ses erreurs. La frontière physique/mathématiques du
\S\ref{sec:motivation} a été traitée comme une contrainte dure sur la formalisation elle-même, non
seulement sur la prose : aucune déclaration Lean de ce développement ne nomme ni ne dépend d'un
objet physique, vérifié par inspection directe et re-vérifié indépendamment par une relecture
contradictoire. Le test même de Tao est pertinent : un résultat que son auteur ne peut ni
expliquer ni défendre, vérifié formellement ou non, n'est pas prêt à être déclaré
complet~\cite{Tao2026AI}. Le contenu mathématique ici --- pourquoi l'auto-dualité effondre le
produit, et pourquoi les deux normalisations du \S\ref{sec:motivation} divergent --- est quelque
chose que l'auteur peut expliquer sans assistance supplémentaire de l'IA ; les termes de preuve
Lean qui le certifient ont été rédigés pour l'essentiel par l'IA et ne sont dignes de confiance
que parce que le noyau, non l'auteur, les vérifie.

\section*{Note sur cette version}

Le présent texte est la version française d'une note rédigée initialement en anglais, dont le
dépôt sur Zenodo est en cours (voir \S\ref{sec:verification} pour l'état de vérification, identique
dans les deux versions). Les deux versions ont le même contenu mathématique ; en cas de
divergence, la version anglaise fait foi pour les énoncés formels, dont les noms Lean ne sont pas
traduits. Cette note n'a pas fait l'objet d'une relecture par les pairs : notre propre relecture
interne a conclu qu'elle ne constituait pas une soumission de revue. Le dépôt HAL vaut archivage
et citation, non validation.

\begin{thebibliography}{99}
\bibitem{Fricke2026} X.~Callens, \textit{A Mathlib-Idiomatic Formalization of the Fricke Involution
and its Operator on Modular Forms}, Zenodo, 2026. \texttt{doi:10.5281/zenodo.22542571}.
\bibitem{EtaQuotients2026} X.~Callens, \textit{Eta-Quotients in Lean 4: Cusp Orders at Every Cusp,
Ligozat's Criterion at Small Level, and a Named Obstruction at General Level}, Zenodo, 2026.
\texttt{doi:10.5281/zenodo.22648098}.
\bibitem{Martin1996} Y.~Martin, \textit{Multiplicative eta-quotients}, Trans.\ Amer.\ Math.\ Soc.\ \textbf{348} (1996), 4825--4856.
\bibitem{PerssonVolpato2015} D.~Persson, R.~Volpato, \textit{Fricke S-duality in CHL models}, JHEP \textbf{12} (2015) 156. \texttt{arXiv:1504.07260}.
\bibitem{Mathlib2020} The mathlib Community, \textit{The Lean mathematical library}, Proceedings of the 9th ACM SIGPLAN International Conference on Certified Programs and Proofs (CPP), 2020.
\bibitem{deMoura2021} L.~de~Moura et al., \textit{The Lean 4 Theorem Prover and Programming Language}, CADE-28, LNCS \textbf{12699}, Springer (2021).
\bibitem{Tao2026AI} T.~Tao, \textit{Mathematics in the age of AI}, essai fondé sur une conférence
au Congrès international des mathématiciens, juillet 2026, soumis aux Actes du CIM 2026.
\texttt{arXiv:2608.16753}.
\bibitem{Leiden2026} \textit{La Déclaration de Leyde sur l'intelligence artificielle et les
mathématiques}, 2 juin 2026, endossée par l'Union mathématique internationale.
\url{https://leidendeclaration.ai/}, \texttt{doi:10.5281/zenodo.20302944}.
\end{thebibliography}

\end{document}
